\section{Discussion}\label{sec:discussion}

% Summarize what we’ve learned
With this paper we aimed to find a common language to categorise different ways to compute model weights, describe the different choices that have to be made together with their implications, and generally highlight their advantages and disadvantages.
We differentiate between two general approaches: ensemble and individual performance weighting.
While in the former weights are computed based on the performance of a linear combination of all models, in the latter, performance is computed for each individual model. As we showed by means of toy examples, a practical advantage of the ensemble performance weighting approach is that it implicitly accounts for the issue of inter-model dependencies since each weight represents the relative contribution of the corresponding model to the weighted average instead of representing its performance as such.
We further considered two additional aspects under the ensemble performance weighting approach, namely the use of multiple diagnostics to measure performance as well as the influence of the choice of the prior distribution over weights.

% multiple diagnostics
When using multiple diagnostics, we showed that applying them jointly so that the weights are tuned to maximize the overall likelihood, i.e.~encompassing all diagnostics, yields better results in terms of the summed RMSE between the weighted average model (using the mean posterior weight vector) and the observed data than applying the diagnostics sequentially and then taking the average of the computed weight vectors. The choice of which diagnostics are most appropriate for measuring performance is a big question on its own \citep[e.g., see][]{glecklerPerformanceMetricsClimate2008} that we have not addressed here. 
Since, in the case of climate projections, we are assessing the performance of physical models, it makes sense to analyze variables that are related to the basic dynamic processes relevant for the target \citep{katzenbergerDevelopingGuidelinesWorking2025}. In any case, it should always be made clear that there is no universal answer and that any probabilistic projections are conditional on the metrics and diagnostic climate variables that were used. Furthermore, when multiple diagnostics are considered, these should be as independent of one another as possible since otherwise we would run into the same problem that we have with co-dependent models in the individual performance weighting approach: good models (diagnostics) would be erroneously overestimated while poor models (diagnostics) would be erroneously underestimated.

Note that in our toy example for the combined diagnostics we standardized the data with respect to the standard deviation of the observations so that the RMSEs are comparable across different diagnostics, i.e.~across different units. This is not necessary as we assume different likelihoods (here different variances) for each diagnostic. When scaling the data by a constant, one should be aware that the variance of the data is scaled by the square of that constant, which also decreases the variance in the weighted averages. This, in turn, decreases the differences between the likelihoods of the observed data under the weighted average models, which, thus, increases the influence of the prior. The predictions of two weighted average models, $m_{*_1}$ and $m_{*_2}$, for example, will differ less when the data is scaled, thus, the respectively used weight vectors, $w_1$ and $w_2$, yield more similar likelihoods. Put differently, when the data is scaled the influence of the likelihood is decreased. A simple possibility to verify the influence of the data is, for instance, to check how sensitive the posterior weights are with respect to different prior distributions. 

Moreover, our example of using ECS as performance metric revealed the importance of prior predictive checks: the prior should be chosen so that the data that is \emph{a priori} generated by the weighted average models is realistic and largely covers the range of plausible data.

% Add paragraph on represented uncertainties
An aspect that we did not consider in more detail in this paper is the representation of the uncertainties. The predicted uncertainty range for our target variable (e.g., future projections of the increase of global mean temperature) based on the computed posterior weight vectors comprises the uncertainty resulting from structural differences between models that is minimized within the ensemble performance weighting approach. What we have not included, but which is possible within a Bayesian formulation, is the uncertainty of the different models themselves (which can be estimated based on multiple model runs) and the uncertainty of the observations.
To include these would be important to get a more realistic estimate of the existing uncertainties.

% Caveats to BMA - when is it not appropriate
A limitation of the ensemble performance weighting approach is that it is restricted to linear combinations of individual models. 
However, in some cases it might be more appropriate to combine models in a non-linear way.
Some models, for instance, perform much better in particular regions of the globe than in others, suggesting that the weights may be modeled to differ across grid points, or regions \citep[e.g.,][]{dasbhowmikPerformancebasedMultimodelCombination2021,thaoCombiningGlobalClimate2022}. 
As the choice of the diagnostics, whether or not it is helpful to compute weights dependent on a region or even on a grid-point level also depends on the objective. If the objective is to compute the likely range of global mean temperatures in the future, it seems more plausible to stick with one weight per model whereas weights on the level of grid cells may be useful when considering likely ranges of temperatures within a certain region only.

A key focus of this work has been to compare the different assumptions and implications of the individual  and the ensemble performance weighting approach. 
Overall, past work \citep{massoudGlobalClimateModel2019,massoudBayesianModelAveraging2020} indicate that ensemble performance weighting provides benefits over individual performance weighting. As just mentioned, there may be cases in which ensemble performance weighting is not appropriate. However, there are many cases in the literature where the final goal is to obtain good performance of the combined ensemble, i.e.~we are interested in knowing the weighted-mean prediction. In those cases, ensemble performance weighting is more straightforwardly connected to the goal and provides an estimate of the uncertainty in the weights which is not accounted for in the individual performance weighting approach. If, on the other hand, the goal is simply to rank models or model versions against each other, then the individual performance weighting approach is straightforward in its application. % and should give the same result as in the ensemble weighting approach.

% Likelihood
%While we used a Gaussian likelihood here, Note though, that BMA is not limited to Gaussian weighting, but any likelihood function can be used.

% Additional considerations possible:
%Independence of observations and gridded model values
%Note Sanderson et al. (2015) section “b. Principal component analysis”, who point to EOFs as a way to retain the right scale of distances. “The use of the EOF prefilter combines fields that are trivially correlated (i.e., adjacent grid cells) into a single mode.”
%Different sources of error and combining them (eg. using a real likelihood function for performance allows explicit representation of error terms…)
%Guidelines

% Final point: BMA probably the best choice in many cases:



